\documentclass[UTF8,fancyhdr,a4paper]{ctexart}

% Title page setup: The following is for titlepage setup 
\title{Linux notes}
\author{李鹏飞}
\date{2020-01-01}

% Margin setup: Set up margin using geometry package
\usepackage[margin = 1in]{geometry}

% Header footer setup: Fancyhdr package sets the header and footer
\usepackage{fancyhdr}
\setlength{\headheight}{40pt}% if you don't have this command, warning for \headheight is too small will appear.
\pagestyle{fancy}% also controls the toc tof tot pages

% 2nd way of setting header
%\pagestyle{myheadings} % this command conflict with the command above
%\markright {李鹏飞\hfill 业余爱好 \hfill}% with the command above, this sets the header

% Fancyhdr details setup
%\fancyhead[l]{我的Logo}% 2nd way of setting header and footer 
\lhead{\includegraphics[width = 1cm]{YinYang.jpg}} % lhead = left head setup
\fancyhead[c]{文档名称}
\fancyhead[r]{文档密级}
\lfoot{12-21-2019}
\cfoot{保密信息，未经授权禁止扩散}
\rfoot{page number}
\usepackage{lastpage}% Use lastpage information
\rfoot{第\thepage 页，共 \pageref{LastPage} 页}
\renewcommand{\footrulewidth}{1pt}% header line width control
\renewcommand{\headrulewidth}{1pt}% footer line width control

% Table of content setup
\usepackage{tocloft}
\tocloftpagestyle{fancy}% to use fancy header and footer, otherwise the tocloft package will overdrive the fancy page setup.
%\usepackage[titles]{tocloft}

\usepackage[xetex]{hyperref}

\begin{document}

%Making title page
\begin{titlepage}
\maketitle
\pagenumbering{gobble}
\centering
\vspace{10cm}
\includegraphics[width = 0.2\textwidth]{YinYang.jpg}\par
\vspace{1cm}
{\huge pengfei.li2017@outlook.com}\par
\vspace{0.5cm}
{\small 版权所有，侵权必究}\par
\vspace{0.5cm}
{\scshape \small All Rights Reserved}
\end{titlepage}

%Decorating toc page
\pagenumbering{roman} % Sets page number in toc, tot, tof
\tableofcontents
\addcontentsline{toc}{section}{目录}%note section here in the second bracket
\clearpage
\listoffigures
\addcontentsline{toc}{section}{插图}
%\addcontentsline{toc}{chapter}{List of Figures}
\clearpage
\listoftables
\addcontentsline{toc}{section}{表格}
%\addcontentsline{toc}{chapter}{List of Tables}
\clearpage

%Writing your document here
\newpage
\pagenumbering{arabic}
\section{Linux common commands}

\subsection{Notes on 20200101}
\begin{enumerate}
\item
check ip addresses of all users in local subnet:\\
sudo nmap -sn 192.168.0.0/24\\
results:\\
Starting Nmap 7.60 ( https://nmap.org ) at 2020-01-01 17:14 CST\\
Nmap scan report for \_ gateway (192.168.0.1)\\
Host is up (0.0043s latency).\\
MAC Address: C8:3A:35:2B:6C:38 (Tenda Technology)\\
Nmap scan report for 192.168.0.100\\
Host is up (0.0040s latency).\\
MAC Address: 94:9F:3E:C3:E1:BD (Sonos)\\
Nmap scan report for 192.168.0.101
Host is up (0.19s latency).\\
MAC Address: 9C:2E:A1:2E:D8:84 (Unknown)\\
Nmap scan report for 192.168.0.103\\
Host is up (-0.053s latency).\\
MAC Address: F4:63:1F:F6:66:43 (Unknown)\\
Nmap scan report for 192.168.0.108\\
Host is up (0.076s latency).\\
MAC Address: BC:83:85:C7:27:73 (Microsoft)\\
Nmap scan report for 192.168.0.109\\
Host is up (0.076s latency).\\
MAC Address: B8:27:EB:26:81:1F (Raspberry Pi Foundation)\\
Nmap scan report for 192.168.0.110\\
Host is up (0.076s latency).\\
MAC Address: A0:51:0B:87:5E:C4 (Unknown)\\
Nmap scan report for pengfei-KLV-WX9 (192.168.0.107)\\
Host is up.\\
Nmap done: 256 IP addresses (8 hosts up) scanned in 4.61 seconds\\




\end{enumerate}

\subsection{Linux remote desktop windows10 home edition}
windows10 home edition 不允许 remote desktop，虽然系统有这个选项询问你是否允许 remote desktop，但是在另一个地方又告诉你 home edition 是没有这个功能的。
以下是解决方案：
\begin{enumerate}
\item 使用 TightVNC:\\
网址：https://www.tightvnc.com/download.php
在 windows10 直接安装 TightVNC，一路安装不用改任何东西，安装完毕。设置require authentication，设置密码。
\item 使用 RDPwrapper 在 github 下又下载:\\
不推荐这个使用方法，CSDN 上很多人反映不可用，而且下载的时候 chrome 会报错。另外还有修改注册表的，不推荐使用。
\item 在 linux 系统下安装 remmina:\\
sudo apt-get install remmina remmina-plugin-*\\
注意尽量需找 apt-get 支持的语句，snap 和 dnf 都不好用，估计是支持其他 linux distro 的。




\end{enumerate}

\subsection{Windows10 remote linux}
需要注意的是我的 raspberryPi 直接允许 VNC 即可，直接通过正常的 windows remote desktop 即可远程登录。但是elementory os 不具有这个功能，于是需要安装。
\begin{enumerate}
\item 在 linux 上直接安装 xrpd：sudo apt-get install xrdp。
\item 直接登录，出现内部错误：这个是由于 VPN 的作用，关掉在开启即可。
\item 登录之后发现闪退：这是由于 linux 没有安装 xfce4。\\ sudo apt-get install xfce4。\par
百度搜索，通过 windows RDP 连接 Linux。
\end{enumerate}


\subsection{Title page 的设置}
使用 begin titlepage 环境，对 title page 进行自定义设置，如图片的插入等;
\subsection{Header 和 Footer 的设置}
这里使用了 fancyhdr package。注意 Lastpage package 的使得 left pagefoot 的页码特特效，第 X 页，共 X 页;
\subsection{toc 或者目录的设置}
使用了 tocloft package。注意这个 package 会覆盖 fancyhdr 的功能，需要命令 tocloftpagestyle（fancy），强行保留 fancyhdr 的功能;\par
hyperef 的使用：注意在 xelatex 编辑环境下。\textbackslash usepackage [xetex] \{ hyperref \},注意 [\space] 中间的选项一定要添加 xetex。\par 
意外发现：\textbackslash[ \space ]\textbackslash ，中间竟然可以输入公式。如下：
\[ x = y^2 \]


\subsection{特殊字符的输入}

\begin{enumerate}
\item
字符输入：一般使用 \textbackslash \space 作为前导。比如 \{ 的输入：Latex 命令为，\textbackslash \{。

\item 
\S \space dollar sign 的使用: 成对使用，用于输入公式，例子如下：
Inline equations input: \$some equations\$, e.g. $\partial T^2/\partial x^2 + \partial T^2/\partial y^2 =\frac{k}{\rho*c_p}*{\partial T/\partial t} $

\item
\textasciitilde 符号的输入有些特殊，必须加入 ascii 关键词：\textbackslash textasciitilde

\end{enumerate}

\subsection{字体的基本设置}



%\addcontentsline{toc}{section}{目录} 这个命令将 toc，tof，tot都加入了目录页中。
\subsection{页码的设置}
使用命令：\textbackslash pagenumbering \{ gobble \}

\newpage
\section{第二个Section}

\subsection{接下来要做的事情}
\begin{enumerate}
\item
URL 的输入
\item
Texworks 的 Linux 安装
\item
Bibliography 的引入
\item
图片的插入及基本设计
\item
表格的插入及基本设置



\end{enumerate}

\newpage
\section{第三个Section}

\newpage
\section{第四个Section}



\end{document}